\documentclass[11pt]{article}

\usepackage{amsmath,amssymb,amsthm,mathtools}
\usepackage[margin=1in]{geometry}

\theoremstyle{plain}
\newtheorem{theorem}{Theorem}
\newtheorem{lemma}[theorem]{Lemma}
\newtheorem{corollary}[theorem]{Corollary}

\theoremstyle{definition}
\newtheorem{definition}[theorem]{Definition}

\newcommand{\assign}{:=}
\newcommand{\CC}{\mathbb{C}}
\newcommand{\RR}{\mathbb{R}}
\newcommand{\ZZ}{\mathbb{Z}}

\title{Autocovariance of the Normalized Zeta-Process}
\author{Stephen Crowley}
\date{April 2026}

\begin{document}
\maketitle

%% ================================================================
\section{Notation and Definitions}
%% ================================================================

Let $\zeta(s)$ denote the Riemann zeta function.
Define the Riemann--Siegel theta function by
\begin{equation}\label{eq:theta-def}
\vartheta(t)
\assign
\operatorname{Im}\log\Gamma\!\bigl(\tfrac{1}{4}+\tfrac{it}{2}\bigr)
- \frac{t}{2}\log\pi.
\end{equation}
Let $t_0$ be the unique point where $\vartheta$ attains its global minimum
on $(0,\infty)$, and set $\tau_0\assign\vartheta(t_0)$.
Let $\sigma:[\tau_0,\infty)\to[t_0,\infty)$ denote the inverse of $\vartheta$
restricted to $[t_0,\infty)$, so that $\vartheta(\sigma(\tau))=\tau$.

The following asymptotic expansions hold as $t\to\infty$:
\begin{align}
\vartheta(t)
&= \frac{t}{2}\log\frac{t}{2\pi} - \frac{t}{2}
   - \frac{\pi}{8} + O\bigl(t^{-1}\bigr),
\label{eq:theta-asymp}\\
\vartheta'(t)
&= \frac{1}{2}\log\frac{t}{2\pi} + O\bigl(t^{-2}\bigr),
\label{eq:thetap-asymp}\\
\vartheta''(t)
&= \frac{1}{2t} + O\bigl(t^{-3}\bigr).
\label{eq:thetapp-asymp}
\end{align}

\begin{definition}\label{def:all}
Define:
\begin{align}
V(T) &\assign \int_0^T
  \bigl|\zeta\bigl(\tfrac{1}{2}+it\bigr)\bigr|^2\,dt,
\label{eq:V}\\
Y(\tau) &\assign
  \frac{\zeta\bigl(\tfrac{1}{2}+i\sigma(\tau)\bigr)}
       {\sqrt{\vartheta'\bigl(\sigma(\tau)\bigr)}},
\label{eq:Y}\\
W(\tau) &\assign \int_{\tau_0}^{\tau} |Y(s)|^2\,ds,
\label{eq:W}\\
X(\tau) &\assign Y(\tau)\sqrt{\frac{\tau}{W(\tau)}}.
\label{eq:X}
\end{align}
\end{definition}

We assume the following results from the companion document:

\smallskip\noindent
\textbf{(F1)}\enspace
$W(\tau) = V(\sigma(\tau)) - V(\sigma(\tau_0))$,
by the substitution $t=\sigma(s)$.

\smallskip\noindent
\textbf{(F2)}\enspace
$W(\tau) = 2\tau\bigl(1+\varepsilon(\tau)\bigr)$
with
$\varepsilon(\tau) = O\!\bigl(\frac{\log\log\tau}{\log\tau}\bigr)$,
obtained by combining (F1) with the Ingham--Atkinson mean-value theorem
$V(T)=T\log\frac{T}{2\pi}+(2\gamma-1)T+E(T)$,
$E(T)=O(T^{1/2}\log^2 T)$,
and the inversion
$\sigma(\tau)=\frac{2\tau}{\log\tau}
\bigl(1+O\bigl(\frac{\log\log\tau}{\log\tau}\bigr)\bigr)$.

\smallskip\noindent
\textbf{(F3)}\enspace
$\displaystyle
\lim_{T\to\infty}\frac{1}{T}\int_{\tau_0}^T |X(\tau)|^2\,d\tau = 1$,
proved via integration by parts on $\int \tau\,d(\log W)$
using only the pointwise asymptotic $\log W(\tau)=\log(2\tau)+o(1)$.

%% ================================================================
\section{Statement of Results}
%% ================================================================

\begin{theorem}\label{thm:cov}
For every fixed $h\in\RR$, the Ces\`aro autocovariance
\begin{equation}\label{eq:Rdef}
R(h)
\assign
\lim_{T\to\infty}\frac{1}{T}
\int_{\tau_0}^T X(\tau+h)\,\overline{X(\tau)}\,d\tau
\end{equation}
exists and equals
\begin{equation}\label{eq:Rvalue}
R(h) = \frac{e^{-ih}\sin h}{h},
\end{equation}
with the convention $R(0)=1$.
\end{theorem}

\begin{corollary}\label{cor:spectral}
The spectral density of $X$ is
$\hat{S}(\omega) = \pi\cdot\mathbf{1}_{[-2,0]}(\omega)$.
\end{corollary}

\begin{proof}[Proof of Corollary~\ref{cor:spectral}]
Using the standard identity
$\int_{-\infty}^{\infty}\frac{\sin h}{h}\,e^{-i\alpha h}\,dh
= \pi\cdot\mathbf{1}_{[-1,1]}(\alpha)$:
\[
\hat{S}(\omega)
= \int_{-\infty}^{\infty} R(h)\,e^{-i\omega h}\,dh
= \int_{-\infty}^{\infty} \frac{\sin h}{h}\,e^{-i(\omega+1)h}\,dh
= \pi\cdot\mathbf{1}_{[-1,1]}(\omega+1)
= \pi\cdot\mathbf{1}_{[-2,0]}(\omega).
\qedhere
\]
\end{proof}

%% ================================================================
\section{Proof of Theorem~\ref{thm:cov}}
%% ================================================================

%% ----------------------------------------------------------------
\subsection*{Separation of the Slowly Varying Prefactor}
%% ----------------------------------------------------------------

From the definitions \eqref{eq:Y} and \eqref{eq:X},
\begin{equation}\label{eq:XbarX}
X(\tau+h)\,\overline{X(\tau)}
= Y(\tau+h)\,\overline{Y(\tau)}
\cdot
\frac{\sqrt{(\tau+h)\,\tau}}
     {\sqrt{W(\tau+h)\,W(\tau)}}.
\end{equation}
Define the prefactor
\begin{equation}\label{eq:alpha-def}
\alpha(\tau,h)
\assign
\frac{\sqrt{(\tau+h)\,\tau}}
     {\sqrt{W(\tau+h)\,W(\tau)}}.
\end{equation}

\begin{lemma}\label{lem:alpha}
For each fixed $h\in\RR$,
\[
\alpha(\tau,h)
= \frac{1}{2} + O\!\left(\frac{\log\log\tau}{\log\tau}\right)
\]
as $\tau\to\infty$.
\end{lemma}

\begin{proof}
By (F2), $W(\tau)=2\tau(1+\varepsilon(\tau))$ and
$W(\tau+h)=2(\tau+h)(1+\varepsilon(\tau+h))$
with $\varepsilon(\tau),\varepsilon(\tau+h)=O(\log\log\tau/\log\tau)$.
Then
\[
\alpha(\tau,h)
= \frac{\sqrt{(\tau+h)\,\tau}}
       {\sqrt{2(\tau+h)(1+\varepsilon(\tau+h))\cdot 2\tau(1+\varepsilon(\tau))}}
= \frac{1}{2}\cdot
  \frac{1}{\sqrt{(1+\varepsilon(\tau+h))(1+\varepsilon(\tau))}}.
\]
Since $(1+u)^{-1/2}=1-u/2+O(u^2)$ for small $u$,
the claim follows.
\end{proof}

Write $\alpha(\tau,h)=\frac{1}{2}+\beta(\tau,h)$
with $\beta(\tau,h)=O(\log\log\tau/\log\tau)$.
Then
\begin{equation}\label{eq:split}
\frac{1}{T}\int_{\tau_0}^T X(\tau+h)\,\overline{X(\tau)}\,d\tau
= \frac{1}{2}\,\mathcal{I}(T,h) + \mathcal{E}(T,h),
\end{equation}
where
\[
\mathcal{I}(T,h)
\assign
\frac{1}{T}\int_{\tau_0}^T Y(\tau+h)\,\overline{Y(\tau)}\,d\tau,
\qquad
\mathcal{E}(T,h)
\assign
\frac{1}{T}\int_{\tau_0}^T
\beta(\tau,h)\,Y(\tau+h)\,\overline{Y(\tau)}\,d\tau.
\]

\begin{lemma}\label{lem:error-vanishes}
$\mathcal{E}(T,h)\to 0$ as $T\to\infty$ for every fixed~$h$.
\end{lemma}

\begin{proof}
By Cauchy--Schwarz,
\[
|\mathcal{E}(T,h)|
\le
\sup_{\tau\ge\tau_0}|\beta(\tau,h)|
\cdot
\left(\frac{1}{T}\int_{\tau_0}^T |Y(\tau+h)|^2\,d\tau\right)^{1/2}
\left(\frac{1}{T}\int_{\tau_0}^T |Y(\tau)|^2\,d\tau\right)^{1/2}.
\]
By (F2), $\frac{1}{T}\int_{\tau_0}^T|Y(\tau)|^2\,d\tau = W(T)/T = 2+o(1)$.
The shifted integral $\frac{1}{T}\int_{\tau_0}^T|Y(\tau+h)|^2\,d\tau
= (W(T+h)-W(h))/T = 2+o(1)$ by the same estimate.
Since $\sup|\beta(\tau,h)| = O(\log\log T/\log T)\to 0$,
the product tends to zero.
\end{proof}

By \eqref{eq:split} and Lemma~\ref{lem:error-vanishes},
to prove Theorem~\ref{thm:cov} it suffices to show
\begin{equation}\label{eq:suffices}
\lim_{T\to\infty}\mathcal{I}(T,h)
= \int_{-2}^{0} e^{i\omega h}\,d\omega
= \frac{2e^{-ih}\sin h}{h},
\end{equation}
since then
$R(h) = \frac{1}{2}\cdot\frac{2e^{-ih}\sin h}{h}
= \frac{e^{-ih}\sin h}{h}$.

%% ----------------------------------------------------------------
\subsection*{Change of Variables to the Height Line}
%% ----------------------------------------------------------------

Set $t=\sigma(\tau)$, so $\tau=\vartheta(t)$ and
$d\tau = \vartheta'(t)\,dt$.
Define the induced height-shift
\begin{equation}\label{eq:delta-def}
\delta_h(t) \assign \sigma\bigl(\vartheta(t)+h\bigr) - t,
\end{equation}
so that $Y(\tau+h)$ evaluated at $\tau=\vartheta(t)$ involves
$\zeta\bigl(\tfrac{1}{2}+i(t+\delta_h(t))\bigr)$.

\begin{lemma}\label{lem:delta}
For each fixed $h$, as $t\to\infty$,
\[
\delta_h(t)
= \frac{h}{\vartheta'(t)}
  + O\!\left(\frac{1}{t(\log t)^3}\right).
\]
In particular,
$\delta_h(t)\,\vartheta'(t) = h + O\bigl(1/(t(\log t)^2)\bigr)$.
\end{lemma}

\begin{proof}
Taylor-expand $\sigma$ at the point $\vartheta(t)$:
\[
\sigma\bigl(\vartheta(t)+h\bigr) - \sigma\bigl(\vartheta(t)\bigr)
= h\,\sigma'\bigl(\vartheta(t)\bigr)
  + \frac{h^2}{2}\,\sigma''(\xi)
\]
for some $\xi$ between $\vartheta(t)$ and $\vartheta(t)+h$.
Since $\sigma'\bigl(\vartheta(t)\bigr) = 1/\vartheta'(t)$
and
\[
\sigma''(s)
= -\frac{\vartheta''\bigl(\sigma(s)\bigr)}
        {\bigl(\vartheta'\bigl(\sigma(s)\bigr)\bigr)^3},
\]
the estimates \eqref{eq:thetap-asymp} and \eqref{eq:thetapp-asymp} give
$\sigma''(\xi) = O\bigl(1/(t(\log t)^3)\bigr)$.
The second claim follows by multiplying through by $\vartheta'(t)$.
\end{proof}

Changing variables $\tau\to t$ in $\mathcal{I}$
and writing $T_*\assign\sigma(T)$ for the upper limit:
\begin{equation}\label{eq:I-height}
\mathcal{I}(T,h)
= \frac{1}{\vartheta(T_*)}
  \int_{t_0}^{T_*}
  \frac{\zeta\bigl(\tfrac{1}{2}+i(t+\delta_h(t))\bigr)\,
        \overline{\zeta\bigl(\tfrac{1}{2}+it\bigr)}}
       {\sqrt{\vartheta'(t+\delta_h(t))\,\vartheta'(t)}}
  \;\vartheta'(t)\,dt.
\end{equation}

\begin{lemma}\label{lem:jacobian}
For fixed $h$,
\[
\frac{\sqrt{\vartheta'(t)}}{\sqrt{\vartheta'(t+\delta_h(t))}}
= 1 + O\!\left(\frac{1}{t\log t}\right).
\]
\end{lemma}

\begin{proof}
Taylor-expand:
$\vartheta'(t+\delta) = \vartheta'(t) + \delta\,\vartheta''(t)
 + O(\delta^2\,\vartheta'''(t))$.
Since $\delta=O(1/\log t)$, $\vartheta''(t)=O(t^{-1})$,
and $\vartheta'''(t)=O(t^{-2})$:
\[
\frac{\vartheta'(t+\delta)}{\vartheta'(t)}
= 1 + O\!\left(\frac{1}{t\log t}\right).
\]
Taking the inverse square root preserves the order.
\end{proof}

Write the ratio in Lemma~\ref{lem:jacobian} as $1+\gamma(t)$
with $\gamma(t)=O(1/(t\log t))$.
After cancelling $\sqrt{\vartheta'(t)}$ from the numerator
and denominator in \eqref{eq:I-height}:
\begin{equation}\label{eq:I-simplified}
\mathcal{I}(T,h)
= \frac{1}{\vartheta(T_*)}
  \int_{t_0}^{T_*}
  \zeta\bigl(\tfrac{1}{2}+i(t+\delta_h(t))\bigr)\,
  \overline{\zeta\bigl(\tfrac{1}{2}+it\bigr)}
  \,dt
  \;+\; \mathcal{J}(T),
\end{equation}
where
\[
\mathcal{J}(T)
= \frac{1}{\vartheta(T_*)}
  \int_{t_0}^{T_*}
  \gamma(t)\,
  \zeta\bigl(\tfrac{1}{2}+i(t+\delta_h(t))\bigr)\,
  \overline{\zeta\bigl(\tfrac{1}{2}+it\bigr)}
  \,dt.
\]
By Cauchy--Schwarz and the mean-value estimate
$\int_{t_0}^{T_*}|\zeta(\tfrac{1}{2}+it)|^2\,dt
= V(T_*) \sim T_*\log T_*$,
together with $\vartheta(T_*)\asymp T_*\log T_*$
and $\sup|\gamma(t)|=O(1/(t_0\log t_0))\to 0$
as the lower integration limit is taken large:
\[
|\mathcal{J}(T)|
\le
\sup_{t\ge t_0}|\gamma(t)|
\cdot
\frac{V(T_*)^{1/2}\cdot V(T_*+|\delta|)^{1/2}}
     {\vartheta(T_*)}
= o(1).
\]
Therefore it suffices to evaluate
\begin{equation}\label{eq:need}
\lim_{T_*\to\infty}
\frac{1}{\vartheta(T_*)}
\int_{t_0}^{T_*}
\zeta\bigl(\tfrac{1}{2}+i(t+\delta_h(t))\bigr)\,
\overline{\zeta\bigl(\tfrac{1}{2}+it\bigr)}
\,dt
= \int_{-2}^{0} e^{i\omega h}\,d\omega.
\end{equation}

%% ----------------------------------------------------------------
\subsection*{Expansion via the Approximate Functional Equation}
%% ----------------------------------------------------------------

For $t$ sufficiently large, the Riemann--Siegel formula gives
\begin{equation}\label{eq:AFE}
\zeta\bigl(\tfrac{1}{2}+it\bigr)
= A(t) + e^{-2i\vartheta(t)}\,\overline{A(t)} + \mathcal{R}(t),
\end{equation}
where
\begin{equation}\label{eq:A-def}
A(t) \assign \sum_{n\le N(t)} n^{-1/2-it},
\qquad
N(t) \assign \Bigl\lfloor\sqrt{\tfrac{t}{2\pi}}\Bigr\rfloor,
\end{equation}
and $\mathcal{R}(t) = O(t^{-1/4})$.
(Here $\overline{A(t)} = \sum_{n\le N(t)} n^{-1/2+it} \eqqcolon B(t)$.)

Abbreviate $\delta=\delta_h(t)$ throughout.
The product
$\zeta\bigl(\tfrac{1}{2}+i(t+\delta)\bigr)\,
 \overline{\zeta\bigl(\tfrac{1}{2}+it\bigr)}$
expands into four main terms plus cross-terms involving $\mathcal{R}$:
\begin{align}
\mathrm{I}   &= A(t+\delta)\,\overline{A(t)},
\label{eq:termI}\\
\mathrm{II}  &= A(t+\delta)\,e^{2i\vartheta(t)}\,B(t),
\label{eq:termII}\\
\mathrm{III} &= e^{-2i\vartheta(t+\delta)}\,B(t+\delta)\,\overline{A(t)},
\label{eq:termIII}\\
\mathrm{IV}  &= e^{-2i\vartheta(t+\delta)+2i\vartheta(t)}\,B(t+\delta)\,\overline{B(t)}.
\label{eq:termIV}
\end{align}

%% ----------------------------------------------------------------
\subsection*{Term~I: Diagonal of the Upper Band}
%% ----------------------------------------------------------------

Expanding \eqref{eq:termI}:
\begin{equation}\label{eq:termI-expand}
A(t+\delta)\,\overline{A(t)}
= \sum_{m,n\le N(t)} \frac{1}{\sqrt{mn}}\,(n/m)^{it}\,m^{-i\delta}.
\end{equation}
(We use $N(t+\delta)=N(t)$ for large $t$ since $\delta\to 0$;
the finitely many terms where $N(t+\delta)\ne N(t)$ contribute
$O(t^{-1/2})$ and are negligible.)

\medskip\noindent
\textit{The diagonal $m=n$.}\enspace
\begin{equation}\label{eq:diagI}
D_{\mathrm{I}}(t)
\assign
\sum_{n\le N(t)} \frac{n^{-i\delta_h(t)}}{n}
= \sum_{n\le N(t)} \frac{e^{-i\delta_h(t)\log n}}{n}.
\end{equation}
By Lemma~\ref{lem:delta},
$\delta_h(t)\log n = \frac{h\log n}{\vartheta'(t)}
+ O\bigl(\frac{\log n}{t(\log t)^3}\bigr)$.
Define the spectral variable
\begin{equation}\label{eq:omega-def}
\omega_n(t) \assign -\frac{\log n}{\vartheta'(t)}.
\end{equation}
For $1\le n\le N(t)$ with $N(t)=\lfloor\sqrt{t/(2\pi)}\rfloor$,
one has $0\le \log n \le \frac{1}{2}\log\frac{t}{2\pi}$,
and by \eqref{eq:thetap-asymp}
$\vartheta'(t) = \frac{1}{2}\log\frac{t}{2\pi} + O(t^{-2})$,
so
\begin{equation}\label{eq:omega-range}
\omega_n(t) \in \bigl[-1-O(t^{-2}),\;0\bigr].
\end{equation}
The phase in \eqref{eq:diagI} is
$-i\delta_h(t)\log n = i\omega_n(t)\,h\,(1+O(1/(t\log^2 t)))$.

To evaluate the sum, apply the Abel summation identity.
Set $L\assign\log N(t) = \frac{1}{2}\log\frac{t}{2\pi}+O(1)$
and $f(u)\assign e^{-iuh}$.
Then:

\begin{lemma}\label{lem:partial-sum}
For a smooth function $f:[0,1]\to\CC$,
\[
\sum_{n\le N}\frac{f\bigl(\frac{\log n}{L}\bigr)}{n}
= L\int_0^1 f(u)\,du + O\!\bigl(\|f\|_\infty + \|f'\|_\infty\bigr),
\]
where $L=\log N$.
\end{lemma}

\begin{proof}
Write $a(x)\assign\sum_{n\le x}1/n = \log x + \gamma + O(1/x)$.
By Abel summation,
\[
\sum_{n\le N}\frac{f\bigl(\frac{\log n}{L}\bigr)}{n}
= f(1)\,a(N) - \int_1^N a(x)\,\frac{d}{dx}\!
  \left[f\!\left(\frac{\log x}{L}\right)\right]dx.
\]
Since $a(x) = \log x + \gamma + O(1/x)$ and
$\frac{d}{dx}f(\frac{\log x}{L}) = \frac{f'(\log x/L)}{xL}$,
the integral is
\[
\int_1^N \frac{(\log x+\gamma)\,f'(\log x/L)}{xL}\,dx + O(1).
\]
Substituting $u=\log x/L$ and using
$\log x = uL$, $dx = xL\,du$:
\[
= \int_0^1 (uL+\gamma)\,f'(u)\,du + O(1).
\]
Integration by parts gives
$\int_0^1 uL\,f'(u)\,du = L\bigl[uf(u)\bigr]_0^1 - L\int_0^1 f(u)\,du$,
and $\int_0^1\gamma\,f'(u)\,du = \gamma(f(1)-f(0))$.
Combining with $f(1)\,a(N) = f(1)(L+\gamma+O(1/N))$:
\[
\sum_{n\le N}\frac{f(\log n/L)}{n}
= L\int_0^1 f(u)\,du + O(\|f\|_\infty + \|f'\|_\infty).
\qedhere
\]
\end{proof}

Apply Lemma~\ref{lem:partial-sum} with $f(u)=e^{-iuh}$
(so $\|f\|_\infty=1$, $\|f'\|_\infty=|h|$)
and note $L = \vartheta'(t) + O(t^{-2})$:
\begin{equation}\label{eq:diagI-value}
D_{\mathrm{I}}(t)
= \vartheta'(t)\int_0^1 e^{-iuh}\,du + O(1)
= \vartheta'(t)\int_{-1}^0 e^{i\omega h}\,d\omega + O(1),
\end{equation}
where the last step is the substitution $\omega=-u$.

\medskip\noindent
\textit{The off-diagonal $m\ne n$.}\enspace
The off-diagonal terms in \eqref{eq:termI-expand} are
\[
\mathrm{OD}_{\mathrm{I}}(t)
= \sum_{\substack{m,n\le N(t)\\m\ne n}}
  \frac{1}{\sqrt{mn}}\,(n/m)^{it}\,m^{-i\delta_h(t)}.
\]
For each pair $(m,n)$ with $m\ne n$,
the factor $(n/m)^{it}=e^{it\log(n/m)}$ oscillates
at frequency $\log(n/m)\ne 0$ with respect to $t$.
The factor $m^{-i\delta_h(t)} = e^{-i\delta_h(t)\log m}$
is slowly varying: since $\delta_h(t)=O(1/\log t)$
and $\log m \le \frac{1}{2}\log t$, its derivative with respect to $t$ is
\[
\frac{d}{dt}\bigl[\delta_h(t)\log m\bigr]
= \delta_h'(t)\log m
= O\!\left(\frac{\log m}{t(\log t)^2}\right)
= O\!\left(\frac{1}{t\log t}\right),
\]
which is $o(1)$ relative to the principal oscillation frequency
$|\log(n/m)|\ge \log(1+1/N(t)) \ge c/\sqrt{t}$.

To control the Ces\`aro average of $\mathrm{OD}_{\mathrm{I}}$,
write $g_{m,n}(t)\assign m^{-i\delta_h(t)}/\sqrt{mn}$
and note $|g_{m,n}(t)|=1/\sqrt{mn}$.
By partial summation in $t$
(integrating $e^{it\log(n/m)}$ and differentiating
the slowly varying factor $g_{m,n}$),
the integral
$\int_{t_0}^{T_*} g_{m,n}(t)\,(n/m)^{it}\,dt$
is bounded by
\[
\frac{2\|g_{m,n}\|_\infty}{|\log(n/m)|}
+ \frac{1}{|\log(n/m)|}
  \int_{t_0}^{T_*}|g_{m,n}'(t)|\,dt
= O\!\left(\frac{1}{\sqrt{mn}\,|\log(n/m)|}\right)
+ O\!\left(\frac{T_*}{\sqrt{mn}\,|\log(n/m)|\,t_0\log t_0}\right).
\]
Summing over all $m\ne n$ with $m,n\le N(T_*)$
and using the Montgomery--Vaughan bound
$\sum_{m\ne n}\frac{1}{\sqrt{mn}\,|\log(n/m)|}
\le C\sum_{n\le N}\frac{1}{n}\cdot\log\log N
= O((\log N)(\log\log N))$
yields
\[
\frac{1}{\vartheta(T_*)}
\int_{t_0}^{T_*}\mathrm{OD}_{\mathrm{I}}(t)\,dt
= O\!\left(\frac{(\log T_*)(\log\log T_*)}{T_*\log T_*}\right)
= o(1).
\]

Combining:
\begin{equation}\label{eq:termI-cesaro}
\frac{1}{\vartheta(T_*)}
\int_{t_0}^{T_*} \mathrm{I}(t)\,dt
= \frac{1}{\vartheta(T_*)}
  \int_{t_0}^{T_*} \vartheta'(t)\,dt
  \cdot\int_{-1}^0 e^{i\omega h}\,d\omega
  + o(1)
= \int_{-1}^0 e^{i\omega h}\,d\omega + o(1),
\end{equation}
since $\int_{t_0}^{T_*}\vartheta'(t)\,dt
= \vartheta(T_*)-\vartheta(t_0)
= \vartheta(T_*)(1+o(1))$.

%% ----------------------------------------------------------------
\subsection*{Term~IV: Diagonal of the Lower Band}
%% ----------------------------------------------------------------

The phase prefactor in \eqref{eq:termIV} is
$e^{-2i\vartheta(t+\delta)+2i\vartheta(t)}$.
By the mean-value theorem and Lemma~\ref{lem:delta}:
\begin{equation}\label{eq:phase-shift}
\vartheta(t+\delta)-\vartheta(t)
= \delta\,\vartheta'(t) + O\bigl(\delta^2\,\vartheta''(t)\bigr)
= h + O\!\left(\frac{1}{t(\log t)^2}\right).
\end{equation}
So the prefactor is
$e^{-2ih}\bigl(1+O(1/(t(\log t)^2))\bigr)$.

The product $B(t+\delta)\,\overline{B(t)}$ expands as
\[
B(t+\delta)\,\overline{B(t)}
= \sum_{m,n\le N(t)} \frac{1}{\sqrt{mn}}\,(m/n)^{it}\,m^{i\delta}.
\]
(Note the signs: $B(t)=\sum n^{-1/2+it}$,
so $\overline{B(t)}=\sum n^{-1/2-it}$.)

The diagonal $m=n$ contributes
\[
D_{\mathrm{IV}}(t)
= e^{-2ih}\sum_{n\le N(t)}\frac{e^{i\delta_h(t)\log n}}{n}
  \cdot\bigl(1+O(1/(t(\log t)^2))\bigr).
\]
The phase per summand is $e^{i\delta_h(t)\log n}
= e^{-i\omega_n(t)\,h\,(1+o(1))}$,
with $\omega_n$ as in \eqref{eq:omega-def} and the sign
reversed relative to Term~I.
Including the prefactor $e^{-2ih}$,
the total phase of the $n$-th term is
$e^{-2ih}\,e^{-i\omega_n h} = e^{i(-\omega_n-2)h}$.

Substituting $\tilde\omega = -\omega_n - 2$,
which maps $\omega_n\in[-1,0]$ to $\tilde\omega\in[-2,-1]$,
and applying Lemma~\ref{lem:partial-sum} in the same manner:
\begin{equation}\label{eq:diagIV-value}
D_{\mathrm{IV}}(t)
= \vartheta'(t)\int_{-2}^{-1} e^{i\omega h}\,d\omega + O(1).
\end{equation}
The off-diagonal terms of $B(t+\delta)\overline{B(t)}$ are controlled
by the same partial-summation and Montgomery--Vaughan argument as
for Term~I (the roles of $m$ and $n$ are interchanged, but the
structure is identical). Therefore:
\begin{equation}\label{eq:termIV-cesaro}
\frac{1}{\vartheta(T_*)}
\int_{t_0}^{T_*}\mathrm{IV}(t)\,dt
= \int_{-2}^{-1} e^{i\omega h}\,d\omega + o(1).
\end{equation}

%% ----------------------------------------------------------------
\subsection*{Terms~II and~III: Rapidly Oscillating Cross-Terms}
%% ----------------------------------------------------------------

Term~II in \eqref{eq:termII} carries the overall phase $e^{2i\vartheta(t)}$,
and Term~III carries $e^{-2i\vartheta(t+\delta)}
= e^{-2i\vartheta(t)}\,e^{-2ih+o(1)}$.
In both cases, the phase $e^{\pm 2i\vartheta(t)}$ has instantaneous
frequency $\pm 2\vartheta'(t)\to\pm\infty$.

The amplitudes (products of Dirichlet polynomials of length $N(t)$)
satisfy
\[
\frac{1}{T_*}\int_{t_0}^{T_*}
|A(t+\delta)||B(t)|\,dt
\le
\left(\frac{V(T_*+|\delta|)}{T_*}\right)^{1/2}
\left(\frac{V(T_*)}{T_*}\right)^{1/2}
= O(\log T_*).
\]

\begin{lemma}\label{lem:riemann-lebesgue}
Let $f:[t_0,\infty)\to\CC$ satisfy
$\frac{1}{T}\int_{t_0}^T |f(t)|\,dt = O((\log T)^A)$
for some $A>0$,
and let $\lambda(t)$ be a real-valued $C^1$ function with
$\lambda'(t)\to\infty$ as $t\to\infty$.
Then
\[
\frac{1}{\lambda(T)}\int_{t_0}^T f(t)\,e^{i\lambda(t)}\,dt \to 0.
\]
\end{lemma}

\begin{proof}
For any $M>0$, split the integral at the point $t_M$
where $\lambda'(t)\ge M$ for all $t\ge t_M$.
The interval $[t_0,t_M]$ contributes $O(1/\lambda(T))=o(1)$.
On $[t_M,T]$, integrate by parts:
\[
\int_{t_M}^T f(t)\,e^{i\lambda(t)}\,dt
= \left[\frac{f(t)\,e^{i\lambda(t)}}{i\lambda'(t)}\right]_{t_M}^T
- \int_{t_M}^T \frac{d}{dt}\!\left(\frac{f(t)}{\lambda'(t)}\right)
  e^{i\lambda(t)}\,dt.
\]
The boundary terms are $O(\|f\|/M)$.
For the integral, the derivative
$\frac{d}{dt}(f/\lambda')$ is integrable
(since $f$ grows at most polynomially and $\lambda'$ grows),
and dividing by $\lambda(T)$ gives a contribution that is $o(1)$
as $T\to\infty$ for any fixed $M$.
Since $M$ is arbitrary, the limit is zero.
\end{proof}

Apply Lemma~\ref{lem:riemann-lebesgue} with
$\lambda(t)=\pm 2\vartheta(t)$
(so $\lambda'(t) = \pm 2\vartheta'(t)\to\pm\infty$)
and $f$ equal to the product of Dirichlet polynomials:
\begin{equation}\label{eq:termII-III-vanish}
\frac{1}{\vartheta(T_*)}
\int_{t_0}^{T_*}\mathrm{II}(t)\,dt \to 0,
\qquad
\frac{1}{\vartheta(T_*)}
\int_{t_0}^{T_*}\mathrm{III}(t)\,dt \to 0.
\end{equation}

%% ----------------------------------------------------------------
\subsection*{Remainder Cross-Terms}
%% ----------------------------------------------------------------

Every cross-term involving $\mathcal{R}(t)=O(t^{-1/4})$ from the
approximate functional equation \eqref{eq:AFE}
produces a product bounded pointwise by
$O(t^{-1/4})\cdot O(t^{1/4+\varepsilon}) = O(t^{\varepsilon})$,
where the second factor uses the convexity bound
$\zeta(\tfrac{1}{2}+it)\ll t^{1/4+\varepsilon}$.
After Ces\`aro averaging:
\begin{equation}\label{eq:remainder}
\frac{1}{\vartheta(T_*)}
\int_{t_0}^{T_*} O(t^{\varepsilon})\,dt
= O\!\left(\frac{T_*^{1+\varepsilon}}{\vartheta(T_*)}\right)
= O\!\left(\frac{T_*^{\varepsilon}}{\log T_*}\right)
= o(1).
\end{equation}

%% ----------------------------------------------------------------
\subsection*{Assembly}
%% ----------------------------------------------------------------

Combining \eqref{eq:termI-cesaro}, \eqref{eq:termIV-cesaro},
\eqref{eq:termII-III-vanish}, and \eqref{eq:remainder}:
\begin{align*}
\lim_{T_*\to\infty}
\frac{1}{\vartheta(T_*)}
\int_{t_0}^{T_*}
\zeta\bigl(\tfrac{1}{2}+i(t+\delta_h(t))\bigr)\,
\overline{\zeta\bigl(\tfrac{1}{2}+it\bigr)}
\,dt
&= \int_{-1}^0 e^{i\omega h}\,d\omega
 + \int_{-2}^{-1} e^{i\omega h}\,d\omega\\
&= \int_{-2}^0 e^{i\omega h}\,d\omega.
\end{align*}
Computing the integral:
\[
\int_{-2}^0 e^{i\omega h}\,d\omega
= \frac{e^{i\omega h}}{ih}\bigg|_{\omega=-2}^{\omega=0}
= \frac{1-e^{-2ih}}{ih}
= \frac{2e^{-ih}\sin h}{h},
\]
where the last equality uses $1-e^{-2ih} = e^{-ih}(e^{ih}-e^{-ih})
= 2ie^{-ih}\sin h$.

By Lemma~\ref{lem:jacobian} and the identity \eqref{eq:I-simplified},
$\mathcal{I}(T,h)$ has the same limit.
By \eqref{eq:split} and Lemma~\ref{lem:error-vanishes}:
\[
R(h)
= \frac{1}{2}\cdot\frac{2e^{-ih}\sin h}{h}
= \frac{e^{-ih}\sin h}{h}.
\]
At $h=0$:
$R(0) = \lim_{h\to 0}\frac{e^{-ih}\sin h}{h} = 1$,
consistent with (F3).
\hfill$\square$

\end{document}
